\chapter{Skin Cancer Excision}

\section*{Introduction \& Clinical Indications}
Excision removes a cutaneous malignancy with a margin of clinically normal tissue and permits histopathologic assessment. It is a common treatment for many basal cell carcinomas, squamous cell carcinomas, and selected melanomas. The correct operation depends on tumor type, size, depth, site, recurrence, perineural or lymphovascular risk, immune status, and whether tissue-sparing or margin-controlled surgery is needed.

Small, well-defined, low-risk lesions may be treated by standard excision. High-risk lesions on the face, ears, hands, genitalia, or other functionally important sites may be better suited to Mohs micrographic surgery or another margin-controlled approach. Melanoma treatment requires a diagnostic biopsy first and a definitive wide local excision based on measured depth; excision should not be planned from visual appearance alone.

\section*{Relevant Surgical Anatomy}
Planning must account for relaxed skin tension lines, cosmetic units, eyelid and lip margins, sensory nerves, tendons, and important vessels. Tumors can extend beyond visible pigment or erythema, and scar tissue from prior biopsy can obscure the original boundaries. The depth of excision varies by tumor and site; preserving uninvolved tissue is important, but an inadequate deep or peripheral margin risks persistence.

\section*{Preoperative Workup \& Preparation}
Confirm pathology whenever feasible and document lesion dimensions, site, biopsy type, prior treatment, and photographs. Examine regional nodes when clinically indicated and review anticoagulants, immunosuppression, wound-healing risks, and allergies. Explain bleeding, infection, nerve injury, incomplete margins, recurrence, contour change, graft or flap failure, and the possible need for re-excision. Select closure only after considering the anticipated defect and the possibility that margins may require further surgery.

\section*{Surgical Technique \& Procedure}
The lesion and planned margin are marked before local anesthetic distorts the field. The specimen is removed with orientation sutures or a diagram so the pathologist can map margins. Dissection proceeds through the intended tissue plane; deep margins must be adequate for the tumor and anatomy. The wound may be closed primarily, left to heal by secondary intention, or reconstructed with a flap or graft. For Mohs surgery, staged mapped sections are examined until margins are clear before reconstruction.

\section*{Complications \& Risks}
Risks include hematoma, infection, dehiscence, hypertrophic scar, altered sensation, injury to motor or sensory nerves, graft or flap necrosis, and functional distortion. Positive margins may require re-excision, margin-controlled surgery, or multidisciplinary treatment. Recurrence risk depends on tumor biology, site, immune status, and completeness of treatment.

\section*{Postoperative Care \& Recovery}
Use pressure and wound-care instructions appropriate to the closure. Review pathology rather than assuming a clinically clear margin is histologically clear. Sutures are removed according to site and tension, and patients receive sun-protection and self-examination advice. Arrange surveillance with dermatology or oncology for high-risk tumors, multiple cancers, melanoma, or immunosuppression.

\section*{Outcomes \& Prognosis}
Early, completely treated nonmelanoma skin cancers usually have excellent local control. Melanoma prognosis is driven mainly by stage and biologic features, not by the size of the skin scar. Reconstruction should restore function and appearance while preserving the ability to monitor for recurrence.

\section*{References}
\begin{enumerate}
\item NCCN. Clinical Practice Guidelines in Oncology: Basal Cell Skin Cancer and Squamous Cell Skin Cancer.
\item NCCN. Clinical Practice Guidelines in Oncology: Cutaneous Melanoma.
\item American Academy of Dermatology. Guidelines of care for the management of cutaneous squamous cell carcinoma and melanoma.
\end{enumerate}
